\section{The setting, spelled out}
\label{sec:setting}

This section fixes notation and, more importantly, states exactly what the existing engine does,
because both algorithms in this document are defined as modifications of it. A reader who already
knows zone graphs can skim to \S\ref{sec:setting:subsumption}, which is where the one fact that
drives everything appears.

\subsection{Networks, valuations, zones}

A \emph{network of timed automata} is a finite set of component automata running in parallel. Its
state has two parts:

\begin{itemize}
  \item a \textbf{discrete state} $q \in \Disc$: the vector of current locations, one per
        component, together with the values of the bounded integer variables. $\Disc$ is finite.
  \item a \textbf{clock valuation} $\nu \in \Val$: one non-negative real per clock.
\end{itemize}

A \emph{concrete state} is a pair $(q,\nu)$. Two things can happen. Time can pass, taking
$(q,\nu)$ to $(q,\nu+d)$ for $d \ge 0$, provided every intermediate valuation satisfies $q$'s
invariant. Or an edge $e$ can fire, which requires $\nu$ to satisfy $e$'s guard, resets the clocks
in $e$'s reset set to zero, and moves to the target discrete state, whose invariant the result must
satisfy.

A \emph{zone} is a set of valuations described by conjunctions of $x_i - x_j \bowtie c$ and
$x_i \bowtie c$. Zones are closed under the operations above and are stored as difference-bound
matrices (DBMs). A \emph{symbolic state} is a pair $(q, \Zn)$ with $\Zn$ a zone; it denotes the set
$\{(q,\nu) : \nu \in \Zn\}$.

The forward image is written $\Post_e$. In the engine it is, in order: intersect with the guard,
apply the resets, intersect with the target invariant, let time pass if the target location is
neither urgent nor committed, intersect with the target invariant again, and finally apply the
abstraction operator $\Extra$ described next. So a stored zone is always \emph{closed under delay}
and always \emph{abstracted}.

\subsection{The abstraction $\Extra$, and why it exists}

Without abstraction the set of reachable zones can be infinite, so the search would not terminate.
The engine applies $\Extra$, an $LU$ extrapolation (\texttt{Extra\_LU+}, and a diagonal-free variant
when the model compares two clocks). $\Extra$ enlarges a zone by forgetting bounds that are beyond
the largest constant the automaton can ever compare that clock against: an upper bound larger than
$U[x]$ becomes $\infty$, and a lower bound smaller than $-L[x]$ is relaxed likewise.

Two properties of $\Extra$ matter later and are stated now so they can be referred to:

\begin{description}
  \item[(A1)] $\Extra$ is \textbf{extensive and monotone}: $\Zn \subseteq \Extra(\Zn)$, and
        $\Zn \subseteq \Zn' \Rightarrow \Extra(\Zn) \subseteq \Extra(\Zn')$. A stored zone is a
        fixpoint of $\Extra$.
  \item[(A2)] $\Extra$ is a \textbf{simulation} abstraction. Every valuation $\Extra$ adds is
        $LU$-simulated by one that was already there. Writing $\nu \simlu \nu'$ for
        ``$\nu'$ simulates $\nu$'', we have
        $\Extra(\Zn) \subseteq \{\nu : \exists \nu' \in \Zn.\ \nu \simlu \nu'\}$.
        This is the standard result of Behrmann, Bouyer, Larsen and Pel\'anek.
\end{description}

(A2) is the single most consequential fact in this document, and \S\ref{sec:soundness} is
essentially an examination of which way it points. Read it carefully: the \emph{added} valuation is
the one being simulated, so the added valuation can do \textbf{at most} what a genuinely reachable
one can do. It cannot do more.

\subsection{Subsumption, and the fact that drives everything}
\label{sec:setting:subsumption}

When a newly generated symbolic state $(q,S)$ is about to be stored and the passed set already
holds $(q,W)$ with $S \subseteq W$ --- or more generally with $S$ $LU$-simulated by $W$ --- the
newcomer is discarded and any edge that pointed at it is redirected to $(q,W)$. This is
\emph{subsumption}, and it is what makes zone-based reachability efficient.

The engine measured here stores, on five of its six benchmark families, \textbf{exactly one zone
per reachable discrete state}, and ties \verifyta{}'s \texttt{States stored} and
\texttt{States explored} counts to the digit at every instance size. Every antichain in its passed
set has exactly one member. We refer to $|\Disc|$, the number of reachable discrete states, as the
\textbf{reachability floor}, and the situation just described as \emph{being on the floor}.

Being on the floor has a sharp consequence that rules out a large class of ideas at once: any
technique whose mechanism is a \emph{comparison between two zones at the same discrete state} ---
a stronger simulation, an exact inclusion test where the abstraction is coarse, a federation or
non-convex union, a wider merge window --- is provably a no-op, because there is never a second
zone to compare against. This is not a prediction; it is what an antichain of size one means.

\subsection{Redirected edges are the problem, not subsumption itself}
\label{sec:setting:subedge}

Call an edge a \textbf{subsumption edge} when the real successor $S$ was strictly smaller than the
stored zone $W$ it was redirected to. The redirection is harmless for reachability: we only ever
ask whether some discrete state is reachable, and $W \supseteq S$ can only make more things
reachable, which is the safe direction for a safety property.

It is fatal for liveness. A cycle in the graph with redirected edges need not correspond to any
cycle the concrete system can execute. The engine has a diagnosed instance of exactly this. On a
railroad-crossing model the search reported the $\lnot q$ cycle
\[
  \texttt{T0.in T1.far \dots Ctl.dn\_1} \;\to\; \texttt{T0.in T1.near \dots Ctl.dn\_2}
  \;\to\; \texttt{T0.in T1.in \dots Ctl.dn\_2} \;\to\; \text{back to the first,}
\]
a train lapping the crossing while \texttt{T0} waits. Each lap costs at least five time units and
\texttt{T0.in} carries the invariant $x_0 \le 15$, so the real system can do three laps, not
infinitely many. What bounds the loop is the \emph{lower} bound on $x_0$, which grows
$0, 5, 10, 15$ and then contradicts the invariant. Subsumption erases precisely that: the real
successor with $x_0 \ge 4$ is a strict subset of the stored zone with $x_0 \ge 0$, so the newcomer
is absorbed and the loop closes on itself.

\begin{remark}[The observation the liveness algorithm is built on]
\label{rem:progress}
What subsumption destroys is a \textbf{monotone progress measure} --- how much of a deadline has
already been consumed. It does not destroy the deadline: the invariant $x_0 \le 15$ is re-applied
at every node, and $\Extra$ dropping it from the stored zone is irrelevant for that reason. So the
lost information is recoverable, and \S\ref{sec:liveness} recovers it with a backward greatest
fixpoint, which is by definition a computation of ``from which valuations can this go on''.
\end{remark}

\subsection{The two queries}

\paragraph{Safety, \texttt{A[] $\varphi$}.} Does every reachable state satisfy $\varphi$?
Decided by exhaustive forward search. All eleven queries in the benchmark suite are of this form.

\paragraph{Leads-to, \texttt{p --> q}.} Equivalent to
$\texttt{A[]}\,(p \Rightarrow \texttt{A<>}\, q)$: does every maximal run from every reachable
$p$-state eventually reach $q$? It is \emph{violated} exactly when some reachable state satisfying
$p \wedge \lnot q$ admits a maximal run that stays in $\lnot q$ forever. A maximal run is either
infinite and time-divergent, or ends where nothing further is possible.

Throughout, $p$ and $q$ are predicates over the \textbf{discrete state only} --- the engine's query
layer rejects clock constraints in them. This restriction is used in
Lemma~\ref{lem:disc-pred} and it is essential; without it the algorithm of \S\ref{sec:liveness}
would need zone splitting at every node.
